\documentclass[10pt,a3paper]{article}
\usepackage[margin=10mm]{geometry}
\usepackage[T1]{fontenc}
\usepackage[utf8]{inputenc}
\usepackage[french]{babel}
\usepackage{lmodern}
\usepackage{microtype}
\makeatletter\def\input@path{{../../../Style/}}\makeatother
\NeedsTeXFormat{LaTeX2e}
\ProvidesPackage{TLPosterCarteMentale}[2026/08/03 Posters et cartes mentales SI]

\RequirePackage{tikz}
\RequirePackage[most]{tcolorbox}
\RequirePackage{xcolor}
\RequirePackage{amsmath,amssymb}
\RequirePackage{graphicx}
\usetikzlibrary{positioning,arrows.meta,calc,shapes.geometric,mindmap,backgrounds,fit}

\definecolor{TLPosterBlue}{RGB}{24,86,141}
\definecolor{TLPosterLight}{RGB}{234,243,250}
\definecolor{TLPosterAccent}{RGB}{232,126,34}
\definecolor{TLPosterGreen}{RGB}{52,152,102}
\definecolor{TLPosterRed}{RGB}{192,57,43}
\definecolor{TLPosterGray}{RGB}{80,90,100}

\newcommand{\TLPosterTitre}[2]{%
  \begin{tcolorbox}[enhanced,boxrule=0pt,arc=0pt,colback=TLPosterBlue,left=6mm,right=6mm,top=4mm,bottom=4mm]
    {\color{white}\bfseries\LARGE #1}\hfill{\color{white}\large #2}
  \end{tcolorbox}%
}

\newtcolorbox{TLPosterBloc}[2][]{
  enhanced,breakable,
  colback=TLPosterLight,
  colframe=TLPosterBlue,
  colbacktitle=TLPosterBlue,
  coltitle=white,
  fonttitle=\bfseries,
  title={#2},
  boxrule=0.8pt,
  arc=2mm,
  left=3mm,right=3mm,top=2mm,bottom=2mm,
  #1
}

\newtcolorbox{TLPosterFormule}[1][]{
  enhanced,
  colback=white,
  colframe=TLPosterAccent,
  boxrule=1pt,
  arc=2mm,
  fontupper=\bfseries,
  halign=center,
  #1
}

\newcommand{\TLPosterMethode}[1]{%
  \begin{tcolorbox}[enhanced,colback=TLPosterGreen!8,colframe=TLPosterGreen,title=Méthode,fonttitle=\bfseries]
  #1
  \end{tcolorbox}%
}

\newcommand{\TLPosterAttention}[1]{%
  \begin{tcolorbox}[enhanced,colback=TLPosterRed!6,colframe=TLPosterRed,title=Point de vigilance,fonttitle=\bfseries]
  #1
  \end{tcolorbox}%
}

\tikzset{
  TLmindroot/.style={concept,concept color=TLPosterBlue,text=white,font=\bfseries\large,minimum size=34mm},
  TLmindlevel1/.style={level 1 concept/.append style={font=\bfseries,minimum size=23mm,sibling angle=45}},
  TLmindlevel2/.style={level 2 concept/.append style={font=\small,minimum size=17mm}},
  TLarrow/.style={-{Stealth[length=3mm]},thick,draw=TLPosterBlue},
  TLbox/.style={rounded corners=2mm,draw=TLPosterBlue,fill=TLPosterLight,align=center,inner sep=3mm}
}

\newenvironment{TLCarteMentale}[1]{%
  \begin{tikzpicture}[mindmap,grow cyclic,concept color=TLPosterBlue,every node/.style={concept},TLmindlevel1,TLmindlevel2]
  \node[TLmindroot]{#1}
}{;
  \end{tikzpicture}
}

\newcommand{\TLBranche}[3][]{child[concept color=#2]{node[#1]{#3}}}

\endinput

\pagestyle{empty}
\renewcommand{\familydefault}{\sfdefault}
\begin{document}
\TLPosterCourseHeader{9}{Machine à courant continu}{Choisir le bon modèle : permanent, transitoire, démarrage ou génératrice}
\vspace{2mm}
\begin{minipage}[t]{0.60\linewidth}
\begin{TLPosterSavoirs}
\begin{itemize}
  \item connaître inducteur, induit, balais--collecteur et rôle du flux magnétique
  \item utiliser $E=K_e\Omega$ et $C_{em}=K_cI$ à flux constant
  \item choisir entre modèle établi et modèle dynamique de l'induit
  \item analyser démarrage, variation de vitesse, défluxage, moteur et génératrice
  \item identifier $K_e$, $R$ et $L$ à partir d'essais adaptés
\end{itemize}
\end{TLPosterSavoirs}
\end{minipage}\hfill
\begin{minipage}[t]{0.37\linewidth}
\TLPosterKey{Quel modèle ?}{Régime établi moteur $\rightarrow U=E+RI$\\Transitoire $\rightarrow U=RI+L\dot I+E$ et PFD\\Démarrage $\rightarrow E=0$\\Génératrice $\rightarrow$ sens de puissance inversé\\Identification $\rightarrow$ vide / rotor bloqué}
\end{minipage}
\vfill
\begin{minipage}[t]{0.49\linewidth}
\begin{TLPosterBloc}[Relations fondamentales]{TLPosterBlue}
À flux constant :
\[
E=K_e\Omega,\qquad C_{em}=K_cI,\qquad K_e=K_c\;\text{numériquement en SI}.
\]
Le \textbf{courant pilote directement le couple}. En régime établi moteur :
\[
U=E+RI,\qquad \Omega=\frac{U-RI}{K_e}.
\]
Le couple utile vaut $C_u=C_{em}-C_p$.
\end{TLPosterBloc}
\begin{TLPosterBloc}[Régime variable : coupler électrique et mécanique]{TLPosterPurple}
Si le courant ou la vitesse évolue, conserver les dynamiques :
\[
U=RI+L\frac{dI}{dt}+K_e\Omega,
\]
\[
J\frac{d\Omega}{dt}=K_cI-C_r-f\Omega.
\]
Le couple résistant est une perturbation mécanique. Ces équations servent à construire le schéma-blocs ou la fonction de transfert du moteur.
\end{TLPosterBloc}
\end{minipage}\hfill
\begin{minipage}[t]{0.49\linewidth}
\begin{TLPosterBloc}[Démarrage et variation de vitesse]{TLPosterGreen}
Au démarrage, $\Omega=0\Rightarrow E=0$ :
\[
I_d=\frac{U}{R},\qquad C_d=K_cI_d.
\]
Le courant doit donc être limité. À flux constant, la vitesse se règle surtout par $U$. Ajouter une résistance série est dissipatif. Au-dessus de la vitesse nominale, le \textbf{défluxage} permet d'augmenter $\Omega$ mais réduit le couple disponible.
\end{TLPosterBloc}
\begin{TLPosterBloc}[Bilan, génératrice et identification]{TLPosterTeal}
Moteur : $UI\rightarrow EI\rightarrow C_{em}\Omega\rightarrow C_u\Omega$ avec pertes Joule $RI^2$ et pertes mécaniques. Génératrice : le flux de puissance s'inverse ; la récupération exige variateur et source réversibles.\\[1mm]
\textbf{À vide :} $U\simeq K_e\Omega$ $\Rightarrow K_e$.\\
\textbf{Rotor bloqué :} $E=0$, réponse RL $\Rightarrow R$ par la valeur finale et $L=R\tau$.
\end{TLPosterBloc}
\end{minipage}
\vfill
\begin{TLPosterMethode}
\begin{enumerate}
  \item Identifier le mode demandé : moteur/génératrice, établi/transitoire, démarrage ou essai d'identification.
  \item Orienter $U$, $I$, $C$ et $\Omega$ puis écrire l'équation électrique adaptée au régime.
  \item Introduire $E=K_e\Omega$ et $C_{em}=K_cI$ ; ne garder $L\dot I$ que si le transitoire l'exige.
  \item Si la vitesse évolue, ajouter l'équation mécanique avec inertie, charge et frottements.
  \item Pour une vitesse imposée, utiliser $\Omega=(U-RI)/K_e$ ; pour un couple imposé, commencer par $I=C_{em}/K_c$.
  \item Terminer par le bilan de puissance et vérifier courant de démarrage, limites nominales, quadrant et réversibilité.
\end{enumerate}
\end{TLPosterMethode}
\vfill
\TLPosterRetenir{MCC à flux constant : le courant fixe le couple et la tension fixe principalement la vitesse ; la première décision consiste donc à choisir le modèle correspondant au régime étudié.}
\end{document}
